\chapter{Происхождение дробного детерминантного барьера из проектных следов}

Предыдущая проверка оставила узкое допустимое окно
\begin{equation}
 0<\nu<\frac{17}{168}
 \label{eq:v6-fractional-window-recall}
\end{equation}
для барьера $-\nu\log\det R$. Теперь выполняется обратный аудит всех
доказанных мер, следов, пфаффианов и малых коэффициентов проекта.

\section{Литературный словарь}

Для конечной алгебры с нормированным следом $\tau$ детерминант
Фугледе--Кадисона имеет логарифм $\tau(\log|A|)$; см.
\path{arXiv:1107.1059}. Поэтому проектор следового веса $w$ передаёт
собственному значению $r$ степень $w\log r$.

Фермионный интеграл зависит от типа полей. Комплексный интеграл даёт
детерминант, а вещественный интеграл --- пфаффиан антисимметричной формы,
квадрат которого равен детерминанту. Для конечных вещественных спектральных
троек этот словарь и его зависимость от KO-степени разобраны в
\path{arXiv:2403.18428}; общая геометрия Pfaffian line обсуждается в
\path{arXiv:1508.04715}.

Существенное ограничение состоит в том, что внутренний KO6-модуль сам по
себе не имеет ненулевого вещественного фермионного интеграла данного типа.
Физический хиральный пфаффиан возникает после произведения с
четырёхмерным пространством, когда суммарная KO-степень равна двум.
Следовательно, ниже строится точный конечномерный коэффициент-кандидат, но
не объявляется уже выведенная полная мера пространства-времени.

\section{Глобальный носитель \texorpdfstring{$M_{300}$}{M300}}

Том V сохранил $M_{300}(\mathbb C)$ как общую алгебру операторов и
единственный носитель нормированного следа
\begin{equation}
 \tau_{300}=\frac1{300}\Tr_{300},
\end{equation}
хотя запретил считать её координатной алгеброй в определяющем
представлении. Одна частичная половина содержит физический сектор
$H_{15}\otimes\mathbb C^3_{\rm fam}$.

Пусть на нём семейное состояние входит как положительный оператор
\begin{equation}
 M_R=I_{15}\otimes R,
 \qquad R>0,\qquad\Tr R=1.
\end{equation}
На двух Real-сопряжённых половинах комплексный модуль детерминанта имеет
тридцать копий семейного логарифма. После нормировки общим следом
получается
\begin{equation}
 \nu_{300}^{\rm det}=\frac{30}{300}=\frac1{10}.
 \label{eq:v6-m300-determinant-weight}
\end{equation}
Если две половины образуют одну антисимметричную форму и интеграл является
пфаффианом, число копий уменьшается вдвое:
\begin{equation}
 \nu_{300}^{\rm Pf}=\frac{15}{300}=\frac1{20}.
 \label{eq:v6-m300-pfaffian-weight}
\end{equation}

Оба коэффициента попадают в окно
\eqref{eq:v6-fractional-window-recall}. Для локального квадратичного
коэффициента получаем точные значения
\begin{equation}
 -\frac{51}{112}+\frac92\frac1{10}
 =-\frac3{560}<0.
 \label{eq:v6-m300-det-residual-curvature}
\end{equation}
\begin{equation}
 -\frac{51}{112}+\frac92\frac1{20}
 =-\frac{129}{560}<0.
 \label{eq:v6-m300-pf-residual-curvature}
\end{equation}
Первый случай почти насыщает границу:
\begin{equation}
 \frac{17}{168}-\frac1{10}=\frac1{840}.
\end{equation}

\begin{proposition}[Условный проход общей меры]
Если полный продукт геометрий индуцирует фермионный детерминант или
пфаффиан семейного оператора $I_{15}\otimes R$ и нормирует его тем же
$\tau_{300}$, то коэффициент барьера фиксирован без подбора и равен
$1/10$ или $1/20$. Оба значения подавляют чистую границу и сохраняют
локальную отрицательную моду в текущем приближении.
\end{proposition}

\section{Локальный тёплицев носитель и расхождение нормировок}

Тёплицев Real-пакет Тома V имеет полный носитель размерности $210$.
Тот же подсчёт даёт
\begin{equation}
 \nu_{210}^{\rm Pf}=\frac{15}{210}=\frac1{14},
 \qquad
 \nu_{210}^{\rm det}=\frac{30}{210}=\frac17.
 \label{eq:v6-m210-fractional-weights}
\end{equation}
Пфаффиановый коэффициент $1/14$ лежит в окне, а полный детерминантный
$1/7$ уже меняет знак локальной кривизны:
\begin{equation}
 -\frac{51}{112}+\frac92\frac17=\frac3{16}>0.
\end{equation}

Это объясняет, почему ранняя догадка о $1/14$ не была чистой нумерологией,
но ещё не делает её выводом. Проект не построил сохраняющее след вложение
локального тёплицева носителя $M_{210}$ в глобальный носитель $M_{300}$ и
не решил, какой след нормирует физический фермионный интеграл. Значения
$1/14$, $1/20$ и $1/10$ отвечают разным корректным контейнерам, а не трём
вариантам одного уже выбранного действия.

\section{Другие найденные совпадения и почему они не переносятся}

Периодический коэффициент $1/24$ Тома II также лежит в допустимом окне.
Но он относится к константной Maxwell--ghost KK-ветви и остаётся условным
результатом выбранной determinant-схемы. Операторной карты от этой
скалярной лапласиановой ветви к $\log\det R$ нет.

Остаточный половинный скалярный детерминант C6 имеет подходящую половинную
степень, но действует на ненулевой скалярной башне
$\mathbb{RP}^3\times S^1$, а не на семейном состоянии. Его перенос в
мостовой сектор изменил бы оператор и спектр.

Наконец, голономный проектор $P_0$ и хиггс-зависимый
$P_\nu(H)$ способны выделить одну нейтринную линию, но после такой
компрессии пфаффиан видит одно матричное собственное значение, а не
$\det R$ всей семейной тройки. Поэтому этот маршрут не создаёт требуемый
барьер до построения полного семейного оператора.

\begin{theorem}[Статус дробного коэффициента]
Проект уже содержит два точных нормировочных механизма, численно попадающих
в допустимое окно: глобальные веса $1/10$, $1/20$ на $M_{300}$ и локальный
пфаффиановый вес $1/14$ на $M_{210}$. Однако ни один ещё не является
родительским барьером, поскольку не построены одновременно:
\begin{enumerate}
 \item полный product-Dirac суммарной KO-степени два;
 \item его антисимметричный семейный блок, зависящий от $R$;
 \item физическая Pfaffian/determinant line;
 \item сохраняющий след переход между локальным и глобальным носителями.
\end{enumerate}
\end{theorem}

Это сильнее простого поиска числа: впервые рабочий диапазон пересекается с
уже доказанными следовыми весами проекта. Но рождение материи остаётся
условным до вычисления одного конкретного фермионного оператора.

\begin{nogocriterion}[Запрет выбора удобного знаменателя]
Нельзя выбрать $210$ или $300$ после сравнения с окном устойчивости,
объявить $1/14$, $1/20$ либо $1/10$ физическим коэффициентом и продолжить.
Следующая проверка обязана построить геометрию произведения и показать, какой
нормированный след возникает в одном функциональном интеграле.
\end{nogocriterion}

\begin{protocol}
\begin{enumerate}
 \item допустимое окно: \textbf{$0<\nu<17/168$};
 \item глобальный determinant-вес $M_{300}$: \textbf{$1/10$};
 \item глобальный Pfaffian-вес $M_{300}$: \textbf{$1/20$};
 \item локальный Pfaffian-вес $M_{210}$: \textbf{$1/14$};
 \item локальный полный determinant-вес $M_{210}$: \textbf{$1/7$};
 \item значения $1/10,1/20,1/14$ внутри окна: \textbf{да};
 \item значение $1/7$ внутри окна: \textbf{нет};
 \item периодический $1/24$ как семейный барьер: \textbf{не типизирован};
 \item внутренний KO6-пфаффиан без пространственного продукта:
 \textbf{не определяет требуемый физический интеграл};
 \item дробный коэффициент как готовый родительский вывод:
 \textbf{ещё нет};
 \item следующая проверка:
 \textbf{product-KO2 пфаффиановый оператор семейного состояния}.
\end{enumerate}
\end{protocol}

Машинный сертификат сохранён в
\path{s2t_v6_fractional_determinant_measure_origin_gate_results.json}.